% Options for packages loaded elsewhere
% Options for packages loaded elsewhere
\PassOptionsToPackage{unicode}{hyperref}
\PassOptionsToPackage{hyphens}{url}
\PassOptionsToPackage{dvipsnames,svgnames,x11names}{xcolor}
%
\documentclass[
  english,
  letterpaper,
  numbers=noendperiod,
  DIV=13]{scrreprt}
\usepackage{xcolor}
\usepackage{amsmath,amssymb}
\setcounter{secnumdepth}{-\maxdimen} % remove section numbering
\usepackage{iftex}
\ifPDFTeX
  \usepackage[T1]{fontenc}
  \usepackage[utf8]{inputenc}
  \usepackage{textcomp} % provide euro and other symbols
\else % if luatex or xetex
  \usepackage{unicode-math} % this also loads fontspec
  \defaultfontfeatures{Scale=MatchLowercase}
  \defaultfontfeatures[\rmfamily]{Ligatures=TeX,Scale=1}
\fi
\usepackage{lmodern}
\ifPDFTeX\else
  % xetex/luatex font selection
\fi
% Use upquote if available, for straight quotes in verbatim environments
\IfFileExists{upquote.sty}{\usepackage{upquote}}{}
\IfFileExists{microtype.sty}{% use microtype if available
  \usepackage[]{microtype}
  \UseMicrotypeSet[protrusion]{basicmath} % disable protrusion for tt fonts
}{}
\makeatletter
\@ifundefined{KOMAClassName}{% if non-KOMA class
  \IfFileExists{parskip.sty}{%
    \usepackage{parskip}
  }{% else
    \setlength{\parindent}{0pt}
    \setlength{\parskip}{6pt plus 2pt minus 1pt}}
}{% if KOMA class
  \KOMAoptions{parskip=half}}
\makeatother
% Make \paragraph and \subparagraph free-standing
\makeatletter
\ifx\paragraph\undefined\else
  \let\oldparagraph\paragraph
  \renewcommand{\paragraph}{
    \@ifstar
      \xxxParagraphStar
      \xxxParagraphNoStar
  }
  \newcommand{\xxxParagraphStar}[1]{\oldparagraph*{#1}\mbox{}}
  \newcommand{\xxxParagraphNoStar}[1]{\oldparagraph{#1}\mbox{}}
\fi
\ifx\subparagraph\undefined\else
  \let\oldsubparagraph\subparagraph
  \renewcommand{\subparagraph}{
    \@ifstar
      \xxxSubParagraphStar
      \xxxSubParagraphNoStar
  }
  \newcommand{\xxxSubParagraphStar}[1]{\oldsubparagraph*{#1}\mbox{}}
  \newcommand{\xxxSubParagraphNoStar}[1]{\oldsubparagraph{#1}\mbox{}}
\fi
\makeatother


\usepackage{longtable,booktabs,array}
\usepackage{calc} % for calculating minipage widths
% Correct order of tables after \paragraph or \subparagraph
\usepackage{etoolbox}
\makeatletter
\patchcmd\longtable{\par}{\if@noskipsec\mbox{}\fi\par}{}{}
\makeatother
% Allow footnotes in longtable head/foot
\IfFileExists{footnotehyper.sty}{\usepackage{footnotehyper}}{\usepackage{footnote}}
\makesavenoteenv{longtable}
\usepackage{graphicx}
\makeatletter
\newsavebox\pandoc@box
\newcommand*\pandocbounded[1]{% scales image to fit in text height/width
  \sbox\pandoc@box{#1}%
  \Gscale@div\@tempa{\textheight}{\dimexpr\ht\pandoc@box+\dp\pandoc@box\relax}%
  \Gscale@div\@tempb{\linewidth}{\wd\pandoc@box}%
  \ifdim\@tempb\p@<\@tempa\p@\let\@tempa\@tempb\fi% select the smaller of both
  \ifdim\@tempa\p@<\p@\scalebox{\@tempa}{\usebox\pandoc@box}%
  \else\usebox{\pandoc@box}%
  \fi%
}
% Set default figure placement to htbp
\def\fps@figure{htbp}
\makeatother


% definitions for citeproc citations
\NewDocumentCommand\citeproctext{}{}
\NewDocumentCommand\citeproc{mm}{%
  \begingroup\def\citeproctext{#2}\cite{#1}\endgroup}
\makeatletter
 % allow citations to break across lines
 \let\@cite@ofmt\@firstofone
 % avoid brackets around text for \cite:
 \def\@biblabel#1{}
 \def\@cite#1#2{{#1\if@tempswa , #2\fi}}
\makeatother
\newlength{\cslhangindent}
\setlength{\cslhangindent}{1.5em}
\newlength{\csllabelwidth}
\setlength{\csllabelwidth}{3em}
\newenvironment{CSLReferences}[2] % #1 hanging-indent, #2 entry-spacing
 {\begin{list}{}{%
  \setlength{\itemindent}{0pt}
  \setlength{\leftmargin}{0pt}
  \setlength{\parsep}{0pt}
  % turn on hanging indent if param 1 is 1
  \ifodd #1
   \setlength{\leftmargin}{\cslhangindent}
   \setlength{\itemindent}{-1\cslhangindent}
  \fi
  % set entry spacing
  \setlength{\itemsep}{#2\baselineskip}}}
 {\end{list}}
\usepackage{calc}
\newcommand{\CSLBlock}[1]{\hfill\break\parbox[t]{\linewidth}{\strut\ignorespaces#1\strut}}
\newcommand{\CSLLeftMargin}[1]{\parbox[t]{\csllabelwidth}{\strut#1\strut}}
\newcommand{\CSLRightInline}[1]{\parbox[t]{\linewidth - \csllabelwidth}{\strut#1\strut}}
\newcommand{\CSLIndent}[1]{\hspace{\cslhangindent}#1}

\ifLuaTeX
\usepackage[bidi=basic]{babel}
\else
\usepackage[bidi=default]{babel}
\fi
% get rid of language-specific shorthands (see #6817):
\let\LanguageShortHands\languageshorthands
\def\languageshorthands#1{}
\ifLuaTeX
  \usepackage[english]{selnolig} % disable illegal ligatures
\fi


\setlength{\emergencystretch}{3em} % prevent overfull lines

\providecommand{\tightlist}{%
  \setlength{\itemsep}{0pt}\setlength{\parskip}{0pt}}



 


\usepackage{grffile}
\KOMAoption{captions}{tableheading}
\makeatletter
\@ifpackageloaded{tcolorbox}{}{\usepackage[skins,breakable]{tcolorbox}}
\@ifpackageloaded{fontawesome5}{}{\usepackage{fontawesome5}}
\definecolor{quarto-callout-color}{HTML}{909090}
\definecolor{quarto-callout-note-color}{HTML}{0758E5}
\definecolor{quarto-callout-important-color}{HTML}{CC1914}
\definecolor{quarto-callout-warning-color}{HTML}{EB9113}
\definecolor{quarto-callout-tip-color}{HTML}{00A047}
\definecolor{quarto-callout-caution-color}{HTML}{FC5300}
\definecolor{quarto-callout-color-frame}{HTML}{acacac}
\definecolor{quarto-callout-note-color-frame}{HTML}{4582ec}
\definecolor{quarto-callout-important-color-frame}{HTML}{d9534f}
\definecolor{quarto-callout-warning-color-frame}{HTML}{f0ad4e}
\definecolor{quarto-callout-tip-color-frame}{HTML}{02b875}
\definecolor{quarto-callout-caution-color-frame}{HTML}{fd7e14}
\makeatother
\makeatletter
\@ifpackageloaded{bookmark}{}{\usepackage{bookmark}}
\makeatother
\makeatletter
\@ifpackageloaded{caption}{}{\usepackage{caption}}
\AtBeginDocument{%
\ifdefined\contentsname
  \renewcommand*\contentsname{Table of contents}
\else
  \newcommand\contentsname{Table of contents}
\fi
\ifdefined\listfigurename
  \renewcommand*\listfigurename{List of Figures}
\else
  \newcommand\listfigurename{List of Figures}
\fi
\ifdefined\listtablename
  \renewcommand*\listtablename{List of Tables}
\else
  \newcommand\listtablename{List of Tables}
\fi
\ifdefined\figurename
  \renewcommand*\figurename{Figure}
\else
  \newcommand\figurename{Figure}
\fi
\ifdefined\tablename
  \renewcommand*\tablename{Table}
\else
  \newcommand\tablename{Table}
\fi
}
\@ifpackageloaded{float}{}{\usepackage{float}}
\floatstyle{ruled}
\@ifundefined{c@chapter}{\newfloat{codelisting}{h}{lop}}{\newfloat{codelisting}{h}{lop}[chapter]}
\floatname{codelisting}{Listing}
\newcommand*\listoflistings{\listof{codelisting}{List of Listings}}
\makeatother
\makeatletter
\makeatother
\makeatletter
\@ifpackageloaded{caption}{}{\usepackage{caption}}
\@ifpackageloaded{subcaption}{}{\usepackage{subcaption}}
\makeatother
\usepackage{bookmark}
\IfFileExists{xurl.sty}{\usepackage{xurl}}{} % add URL line breaks if available
\urlstyle{same}
\hypersetup{
  pdftitle={Insights of Particle Image Velocimetry},
  pdfauthor={JF Krawczynski},
  pdflang={en},
  colorlinks=true,
  linkcolor={blue},
  filecolor={Maroon},
  citecolor={Blue},
  urlcolor={Blue},
  pdfcreator={LaTeX via pandoc}}


\title{Insights of Particle Image Velocimetry}
\author{JF Krawczynski}
\date{2025-10-03}
\begin{document}
\maketitle

\renewcommand*\contentsname{Table of contents}
{
\hypersetup{linkcolor=}
\setcounter{tocdepth}{2}
\tableofcontents
}

\bookmarksetup{startatroot}

\chapter*{About}\label{about}
\addcontentsline{toc}{chapter}{About}

\markboth{About}{About}

This site hosts the course materials for fluid mechanics measurements by
PIV for the course unit \textbf{UM5MEE12 - Experimental Methods and Data
Processing} taught in the
\href{https://masters-sdi.sorbonne-universite.fr/la-mention-mecanique}{Master
of Engineering Sciences} at the
\href{https://sciences.sorbonne-universite.fr/faculte/ufr-instituts-observatoires-ecoles/ufr-dingenierie}{Faculty
of Engineering} of \href{https://www.sorbonne-universite.fr/}{Sorbonne
University}. The aim is to understand the basic concepts inherent to
measuring velocity fields in fluid flows.

\begin{tcolorbox}[enhanced jigsaw, titlerule=0mm, colframe=quarto-callout-note-color-frame, bottomtitle=1mm, leftrule=.75mm, opacityback=0, bottomrule=.15mm, arc=.35mm, colbacktitle=quarto-callout-note-color!10!white, coltitle=black, breakable, colback=white, opacitybacktitle=0.6, toptitle=1mm, rightrule=.15mm, left=2mm, toprule=.15mm, title=\textcolor{quarto-callout-note-color}{\faInfo}\hspace{0.5em}{Note}]

This course is partly adapted from the
\href{https://github.com/OpenPIV/openpiv-python}{documentation} of
\href{https://doi.org/10.5281/zenodo.4409178}{OpenPIV}, licensed under
the GNU General Public License v2.0.

\end{tcolorbox}

The course content is divided into two parts:

\begin{itemize}
\tightlist
\item
  The theoretical aspects of the course, including:

  \begin{itemize}
  \tightlist
  \item
    the fundamental physical considerations related to the tracer
    particles used to measure the flow
  \item
    the algorithms developed to obtain the velocity vector field from
    the analysis of particle images
  \end{itemize}
\item
  The practical aspects:

  \begin{itemize}
  \tightlist
  \item
    a numerical part to better understand the evaluation methods using
    PIV. The images processed will be generated synthetically
  \item
    an experimental part to understand the trade-offs required to obtain
    high-quality images for processing
  \end{itemize}
\end{itemize}

In this course, we will mainly use the following tools, with which
students should be familiar: - \href{https://www.python.org/}{Python} -
\href{https://numpy.org}{Numpy} -
\href{https://matplotlib.org}{Matplotlib}

\bookmarksetup{startatroot}

\chapter*{Assessment}\label{assessment}
\addcontentsline{toc}{chapter}{Assessment}

\markboth{Assessment}{Assessment}

Assessment will be based on a short lab report (for the numerical part)
and a project report focusing on the analysis of experimental data
obtained during the lab sessions.

\bookmarksetup{startatroot}

\chapter*{Schedule}\label{schedule}
\addcontentsline{toc}{chapter}{Schedule}

\markboth{Schedule}{Schedule}

Here is the planned organization of the sessions:

\begin{longtable}[]{@{}
  >{\raggedleft\arraybackslash}p{(\linewidth - 4\tabcolsep) * \real{0.3125}}
  >{\centering\arraybackslash}p{(\linewidth - 4\tabcolsep) * \real{0.3125}}
  >{\raggedright\arraybackslash}p{(\linewidth - 4\tabcolsep) * \real{0.3750}}@{}}
\toprule\noalign{}
\begin{minipage}[b]{\linewidth}\raggedleft
Week
\end{minipage} & \begin{minipage}[b]{\linewidth}\centering
Date
\end{minipage} & \begin{minipage}[b]{\linewidth}\raggedright
Topic
\end{minipage} \\
\midrule\noalign{}
\endhead
\bottomrule\noalign{}
\endlastfoot
1 & 09/24/2025 & Introduction to velocity field measurement by PIV,
Lecture 1 \\
2 & 10/01/2025 & Algorithms and experimental specifics, Lab 1 + Lab 2 \\
3 & 10/08/2025 & Algorithms and experimental specifics, Lab 1 + Lab 2 \\
\end{longtable}

\bookmarksetup{startatroot}

\chapter*{Recommended Readings and
Media}\label{recommended-readings-and-media}
\addcontentsline{toc}{chapter}{Recommended Readings and Media}

\markboth{Recommended Readings and Media}{Recommended Readings and
Media}

Below is a list of useful resources for your project, organized by
topic.

\subsubsection*{Python for Scientific
Programming}\label{python-for-scientific-programming}
\addcontentsline{toc}{subsubsection}{Python for Scientific Programming}

\begin{itemize}
\tightlist
\item
  \href{https://zestedesavoir.com/tutoriels/2514/un-zeste-de-python/}{Un
  Zeste de Python}
\item
  \href{https://zestedesavoir.com/tutoriels/4139/les-bases-de-numpy-et-matplotlib/}{Les
  bases de numpy et matplotlib}
\item
  \href{https://www.labri.fr/perso/nrougier/from-python-to-numpy/}{From
  Python to Numpy}
\item
  \href{https://inria.hal.science/hal-03427242v1/file/scientific-visualization-python-matplotlib.pdf}{Scientific
  Visualization: Python + Matplotlib} (N. P. Rougier 2021)
\end{itemize}

\subsubsection*{Scientific outputs: figures and
reports}\label{scientific-outputs-figures-and-reports}
\addcontentsline{toc}{subsubsection}{Scientific outputs: figures and
reports}

\begin{itemize}
\tightlist
\item
  \href{https://journals.plos.org/ploscompbiol/article/file?id=10.1371/journal.pcbi.1003833&type=printable}{Ten
  Simple Rules for Better Figures} (M. D. Rougier Nicolas P. 2014)
\item
  \href{https://fr.wikibooks.org/wiki/LaTeX}{Wikibook LaTeX}
\item
  \href{https://typst.app/}{Typst}
\end{itemize}

\part{Theory}

\chapter{Introduction}\label{introduction}

This is a book created from markdown and executable code.

See Knuth (1984) for additional discussion of literate programming.

\part{Practice}

\chapter{Summary}\label{summary}

In summary, this book has no content whatsoever.

\bookmarksetup{startatroot}

\chapter*{References}\label{references}
\addcontentsline{toc}{chapter}{References}

\markboth{References}{References}

\phantomsection\label{refs}
\begin{CSLReferences}{1}{0}
\bibitem[\citeproctext]{ref-knuth84}
Knuth, Donald E. 1984. {``Literate Programming.''} \emph{Comput. J.} 27
(2): 97--111. \url{https://doi.org/10.1093/comjnl/27.2.97}.

\bibitem[\citeproctext]{ref-rougierTenSimpleRules2014}
Rougier, Michael Droettboom, Nicolas P. 2014. {``Ten Simple Rules for
Better Figures.''} \emph{PLOS Computational Biology} 10 (9).
\url{https://doi.org/10.1371/journal.pcbi.1003833}.

\bibitem[\citeproctext]{ref-rougierScientificVisualizationPython2021}
Rougier, Nicolas P. 2021. \emph{Scientific Visualization: Python +
Matplotlib}. CRC Press.

\end{CSLReferences}




\end{document}
